\subsection*{Реалізація ключових алгоритмів}

\subsection{Геометричний контролер на SE(3)}

Клас \texttt{CommandedNonlinearController} (файл \texttt{swarm/controller.py}) є підкласом \texttt{pegasus.NonlinearController}~\cite{Pegasus2024}. Він не змінює внутрішньої математики контролера, а лише замінює джерело референсного стану: замість CSV-таблиці викликає \texttt{TrajectoryFollower.advance(dt)} і передає у батьківський \texttt{update} шестиелементний вектор:
\[
\bigl(p_{\text{ref}},\; v_{\text{ref}},\; a_{\text{ref}},\; j_{\text{ref}},\; \psi_{\text{ref}},\; \dot{\psi}_{\text{ref}}\bigr)
\]

Внутрішній контур \texttt{NonlinearController} обчислює:

\textbf{Бажана сила:}
\[
F_{\text{des}} = -K_p e_p - K_d e_v - K_i \textstyle\int e_p\,dt + mg\mathbf{e}_3 + m a_{\text{ref}}
\]
де $e_p = p - p_{\text{ref}}$, $e_v = v - v_{\text{ref}}$.

\textbf{Тяга і бажана орієнтація:}
\[
u_1 = F_{\text{des}}^\top R\,\mathbf{e}_3, \qquad
Z^b_d = \frac{F_{\text{des}}}{\|F_{\text{des}}\|}, \qquad
R_d = [X^b_d, Y^b_d, Z^b_d]
\]

\textbf{Помилка орієнтації та моменти:}
\[
e_R = \tfrac{1}{2}(R_d^\top R - R^\top R_d)^\vee, \qquad
\tau = -K_R e_R - K_w e_w
\]

\textbf{Алокація} $(u_1, \tau) \to (\omega_1,\ldots,\omega_4)$ — штатна матриця Pegasus \cite{Pegasus2024}.

Усі коефіцієнти ($K_p, K_d, K_i, K_R, K_w$) і параметри ($m, g$) — за замовчуваннями Pegasus.

\subsection{Квінтична поліноміальна траєкторія}

Для одного сегмента від точки $p_0$ до $p_1$ будуємо поліном 5-го ступеня в кожній осі з граничними умовами $\{p, v, a\}$ на початку і кінці (рівняння 9.13–9.16 з~\cite{Lynch2017}):
\[
p(t) = c_0 + c_1 t + c_2 t^2 + c_3 t^3 + c_4 t^4 + c_5 t^5
\]
Тривалість сегмента:
\[
T = \max\!\left(T_{\min},\; \frac{\|p_1 - p_0\|}{v_{\text{cruise}}}\right), \quad T_{\min} = 2{,}0\;\text{с}
\]

Для одиночних команд (\texttt{takeoff}, \texttt{goto}, \texttt{land}): $v_0 = v_1 = 0$, $a_0 = a_1 = 0$.

Для ланцюжка waypoints (\texttt{scan}): на стиках $v_i$ є вектором бісектриси, нормованим до $v_{\text{cruise}}$ — це забезпечує C¹-неперервність при прагматичному відмові від повного QP-minimum-snap~\cite{Mellinger2011}.

\subsection{Boustrophedon Coverage (BCD)}

Функція \texttt{boustrophedon\_path} у \texttt{swarm/coverage.py} реалізує примітив~\cite{Choset2000}:

\begin{lstlisting}[language=Python, caption={Генерація BCD-waypoints для коридора}]
def boustrophedon_path(
    rect: Rectangle,
    altitude_m: float,
    sweep_spacing_m: float,
    sweep_axis: str = "x",
) -> list[np.ndarray]:
    # Generate stripe centroids along the primary axis
    xs = np.arange(rect.x_min, rect.x_max + 1e-6, sweep_spacing_m)
    waypoints = []
    for i, x in enumerate(xs):
        if i % 2 == 0:  # forward sweep
            waypoints += [(x, rect.y_min, altitude_m),
                          (x, rect.y_max, altitude_m)]
        else:            # reverse sweep (boustrophedon)
            waypoints += [(x, rect.y_max, altitude_m),
                          (x, rect.y_min, altitude_m)]
    return [np.array(wp) for wp in waypoints]
\end{lstlisting}

Для вузьких коридорів ($\Delta x \approx 2{,}5$~м): \texttt{x\_min}~= центр коридора, \texttt{x\_max}~= \texttt{x\_min + 0.05}, \texttt{sweep\_spacing\_m = 10.0} — генерується рівно 2 waypoints (один прямолінійний прохід).

\subsection{Орбіта навколо точки інтересу}

Функція \texttt{build\_orbit} у \texttt{swarm/trajectory.py} генерує $N=16$ рівновіддалених точок на колі:
\[
\theta_k = \frac{2\pi k}{N}, \quad
p_k = \bigl(x_c + r\cos\theta_k,\; y_c + r\sin\theta_k,\; z_c\bigr),\quad k=0,\ldots,N
\]
Yaw у кожній точці спрямований до центра (inward-facing camera):
$\psi_k = \pi + \theta_k$. Замкнений контур: $p_N = p_0$.

\subsection{Естафетна передача місії}

Модуль \texttt{swarm/relay.py} містить два класи.

\textbf{DroneStateTracker} підписується на \texttt{swarm/state/\{drone\_id\}} і в callback оновлює структуру \texttt{WorldPose(x, y, z, battery\_pct, mode)}:

\begin{lstlisting}[language=Python, caption={Thread-safe читання позиції}]
class DroneStateTracker:
    def snapshot(self) -> WorldPose:
        with self._lock:
            return dataclasses.replace(self._pose)
\end{lstlisting}

\textbf{RelayMission} виконує послідовність:

\begin{enumerate}[itemsep=2pt]
  \item Дрон~A: \texttt{takeoff} → \texttt{goto}(вхід коридора) → \texttt{goto}(кінець коридора).
  \item Паралельно — пулінг \texttt{DroneStateTracker} кожні 0,5~с. При \texttt{battery\_pct < threshold}:
  \begin{enumerate}
    \item \texttt{stop}(A), \texttt{sleep(1.0)} — дрон стабілізується.
    \item \texttt{handoff\_pose = tracker.snapshot()} — \emph{поточні світові координати}.
    \item \texttt{land}(A).
    \item Дрон~B: \texttt{takeoff} → \texttt{goto}(\texttt{handoff\_pose.xyz}) → \texttt{goto}(кінець коридора) → \texttt{land}.
  \end{enumerate}
\end{enumerate}

Ключовий момент: позиція перехоплення — це не абстрактний ID, а реальні XYZ-координати зі MQTT-повідомлення стану дрона~A.

\begin{lstlisting}[language=Python, caption={Перехоплення поточних координат}]
handoff = tracker_a.snapshot()
print(f"Handoff world pose:  "
      f"X={handoff.x:.2f}  Y={handoff.y:.2f}  Z={handoff.z:.2f}")

cmd = GoToCmd(x=handoff.x, y=handoff.y, z=handoff.z,
              speed_mps=relay_speed)
publish_command("uav-02", cmd)
\end{lstlisting}

\subsection{Симуляція заряду батареї}

Заряд симулюється лінійно:
\[
\text{battery}(t+\Delta t) = \text{battery}(t) - \text{drain\_rate} \times \Delta t
\]
де \texttt{drain\_rate} (в~\%/с) задається через \texttt{SIM\_BATTERY\_DRAIN\_RATE\_PPS} у \texttt{.env}.

Реалістичне значення: $0{,}05$~\%/с ($\approx 33$~хв до повного розряду).
Для демонстрації: $2{,}0$~\%/с (поріг 60\% досягається через $\approx 20$~с польоту).

\subsection{Хмарний стек моніторингу}

\begin{itemize}
  \item \textbf{Telegraf}: MQTT Consumer плагін підписується на \texttt{swarm/telemetry/+} та записує InfluxDB Line Protocol безпосередньо у bucket \texttt{telemetry}.
  \item \textbf{InfluxDB v2}: зберігає часові ряди; Flux-запити Grafana звертаються до measurement \texttt{drone\_telemetry} (поля: \texttt{x}, \texttt{y}, \texttt{z}, \texttt{speed}, \texttt{battery\_pct}, \texttt{coverage\_pct}).
  \item \textbf{Grafana}: 8 панелей; оновлення кожні 2~с; threshold для \texttt{battery\_pct} < 30\% — червоний колір.
\end{itemize}
